% Part: many-valued-logic
% Chapter: syntax-and-semantics
% Section: connectives

\documentclass[../../../include/open-logic-section]{subfiles}

\begin{document}

\olfileid{mvl}{syn}{con}

\olsection{زباناں تے رابطے}

کلاسیکی قضیاتی منطق تے کئی ہور منطقاں وچ \emph{قضیاتی مستقل}
تے \emph{رابطے} اک مقرر گنتی وچ ورتے جاندے نیں۔ مثلاً اسیں
ہِٹھلے نشان مُڈھلے نشان وجوں ورتدے آں۔
\begin{enumerate}
\tagitem{prvFalse}{!!{falsity} لئی قضیاتی مستقل~$\lfalse$۔}{}
\tagitem{prvTrue}{!!{truth} لئی قضیاتی مستقل~$\ltrue$۔}{}
\item منطقی رابطے۔
  \startycommalist
  \iftag{prvNot}{\ycomma $\lnot$ (نفی)}{}%
  \iftag{prvAnd}{\ycomma $\land$ (عطف)}{}%
  \iftag{prvOr}{\ycomma $\lor$ (فصل)}{}%
  \iftag{prvIf}{\ycomma $\lif$ (!!{conditional})}{}%
  \iftag{prvIff}{\ycomma $\liff$ (!!{biconditional})}{}%
\end{enumerate}
\iftag{defNot,defOr,defAnd,defIf,defIff,defTrue,defFalse,defEx,defAll}{%
اُپرلے مُڈھلے رابطیاں دے نال اسیں اوہ نشان وی ورتدے آں
جیہڑے مختصر لکھتاں وجوں متعین ہوندے نیں۔ مثلاً
\startycommalist
  \iftag{defNot}{\ycomma $\lnot$ (نفی)}{}%
  \iftag{defAnd}{\ycomma $\land$ (عطف)}{}%
  \iftag{defOr}{\ycomma $\lor$ (فصل)}{}%
  \iftag{defIf}{\ycomma $\lif$ (!!{conditional})}{}%
  \iftag{defIff}{\ycomma $\liff$ (!!{biconditional})}{}%
  \iftag{defFalse}{\ycomma $\lfalse$ (!!{falsity})}{}%
  \iftag{defTrue}{\ycomma $\ltrue$ (!!{truth})}۔}{}

کئی قدراں والیاں منطقاں وچ وی اوہی رابطے ورتے جاندے نیں۔
پر اکثر اکّو منطق وچ، مثلاً عطف دے وکھ وکھ روپ رکھنا کم
آؤندا اے۔ روپاں نوں وکھ رکھن لئی وکھرے نشان چاہیدے نیں۔
کجھ کئی قدراں والیاں منطقاں وچ ایسے رابطے وی ہوندے نیں
جنہاں دا کلاسیکی منطق وچ کوئی ہم ارز نہیں۔ ایس لئی اسیں
عام نالوں ذرا ودھ کھلھ نال تعریف کراں گے۔

\begin{defn}
  \emph{قضیاتی زبان} \emph{رابطیاں} دے سیٹ $\Lang L$ توں
  بندی اے۔ ہر رابطے $\star$ دی اک \emph{مقاماں دی گنتی} اے۔
  جیہدے مقاماں دی گنتی~$n$ اے، اوہنوں \emph{$n$-مقامی} رابطہ
  آکھدے آں۔ مقاماں دی گنتی~$0$ ہووے تاں اوہ رابطے
  \emph{مستقل} وی اکھواؤندے نیں۔ مقاماں دی گنتی~$1$ ہووے
  تاں \emph{اک مقامی}، تے مقاماں دی گنتی~$2$ ہووے تاں
  \emph{دو مقامی} اکھواؤندے نیں۔
\end{defn}

\begin{ex}
  قضیاتی منطق دی معیاری زبان $\Lang L_0$ وچ ہِٹھلے رابطے نیں۔
  نال اوہناں دے مقاماں دی گنتی دِتّی اے۔
  $\lfalse$~($0$)،
  $\lnot$~($1$)،
  $\land$~($2$)،
  $\lor$~($2$)،
  $\lif$~($2$)۔ اسیں جِنہاں منطقاں نوں ویکھاں گے، اوہناں
  وچوں بہتیاں وچ ایہو زبان ورتی جاندی اے۔ رواج مطابق کجھ
  منطقاں وچ کجھ رابطیاں دے نشان وکھرے نیں۔ مثلاً حاصلِ ضرب
  والی منطق وچ عطف دا نشان اکثر $\odot$ ہوندا اے~$\land$
  دی تھاں۔ کدے نواں عامل شامل کرنا سوکھا پیندا اے، جیویں
  تعیّن دا عامل $\triangle$ ($1$-مقامی)۔
\end{ex}

\end{document}
